% Limitations. Written before results; every number is checked by tests/test_paper_limitations.py.
\section{Limitations}
\label{sec:limitations}

\paragraph{Pilot scale.}
The protocol planned a main experiment with about $5{,}000$ training root graphs. It was not
run: all adapters were trained on $512$ root graphs for $3$ epochs, and the locked test has
$200$ root graphs per split ($100$ for the combined split). Whether the ordering of the arms
holds with ten times the training data is not measured here.

\paragraph{Synthetic, clean drawings.}
Every image is a procedurally drawn $512 \times 512$ directed graph with readable node
spacing and no noise, handwriting, scanning artefacts or domain text. Nothing in this study
measures how a model reads real diagrams, and no claim is made about them.

\paragraph{One base model, one adapter type, one edit budget.}
All arms fine-tune the same 2B base model (Qwen3.5-2B) with LoRA under one setting, and the
repair task allows at most one edge change. A different base model, full fine-tuning, or
multi-edit repair could behave differently; these were not run.

\paragraph{Twins protect only the main tasks.}
Answer-flipping twins cap an image-blind predictor at $0.50$ on the three main tasks. The
secondary tasks have no twins: on the pilot development split a per-instruction lookup
already reaches $0.895$ on edge existence. Accuracy on secondary tasks therefore cannot
separate use of the image from guessing, and we do not use them for the main comparison.

\paragraph{Unequal supervision per arm.}
Arms are matched on samples, root graphs and images per epoch, not on answer tokens. On the
pilot training split the repair arm D has $16{,}384$ target tokens per epoch, the
question-answering arm B $5{,}120$ and the structural arm C $28{,}592$
(Section~\ref{sec:setup}). A difference between arms can partly reflect how much answer
text each one supervises.

\paragraph{Generator coverage and selection.}
Small graphs run out of distinct topologies: about $100$ five-node classes for
\texttt{layered\_dag}, $255$ for \texttt{tree\_plus} and $422$ for \texttt{sparse\_random}.
The planned main build therefore starts at seven nodes; the pilot graphs used here have 5 to 10. Layout rejection also prefers sparser
graphs at large node counts, so the size-OOD split differs from the training graphs in
density as well as size; we report edges per root and a density-matched slice
(Section~\ref{sec:analysis}) rather than claiming a clean size effect. Layout is not held
out in the style-OOD split.

\paragraph{Controls not run.}
The protocol lists a question-only (blank image) control and a wrong-image control. Neither
was run on any model; only model-free text probes exist (the ceilings above). That the arms
use the image therefore rests on the twin design (an image-blind predictor cannot exceed
$0.50$ on the main tasks), not on an ablation of the image. The oracle condition (graph given as text) was evaluated on $100$
development roots for arms B and D and not on the locked test. Slices cover graph family,
size, layout and style; there is no slice by path length or by edge crossings.

\paragraph{Reproducibility across machines.}
Graphs, labels and instructions rebuild to the same content digest on every machine we
tried, but PNG bytes differ with font rasterisation. Training runs used cloud T4 GPUs (Colab, Kaggle)
whose library versions and hardware can differ between providers; we record the
environment of each run and compare runs only within one environment unless stated.
Within one environment, training is repeatable: retraining arm D with its recorded seed from a fresh
clone of the repository on a Kaggle T4 gave byte-identical development predictions to the
recorded Kaggle run (\texttt{scripts/repro\_compare.py}, \texttt{reports/repro.md}).

\paragraph{Small budget.}
The study runs on single T4 GPUs with no extra spending. This limits the number of seeds and the size of the
training sets, so intervals are wide and small effects may go undetected.
